\begin{table}[H]
    \centering
    \caption{Desempenho dos classificadores no conjunto de teste.}
    \label{tab:comparacao_classificadores_classicos}
    \footnotesize
    \begin{tabular}{lccc}
        \hline
        \textbf{Classificador} & \textbf{Acurácia (\%)} &
        \textbf{Acurácia balanceada (\%)} & \textbf{F1 macro (\%)} \\
        \hline
        Random Forest & \textbf{92,37} & \textbf{92,20} & \textbf{92,32} \\
        SVM linear & 91,94 & 91,65 & 91,83 \\
        Extra Trees & 90,83 & 90,73 & 90,80 \\
        MLP & 87,45 & 87,56 & 87,45 \\
        \hline
    \end{tabular}
    \fonte{Elaborado pelo autor (2026).}
    \vspace{0.15cm}
    \parbox{0.96\textwidth}{\footnotesize\textit{Nota:} todos os modelos foram
    avaliados nas mesmas 22.115 leituras de teste, pertencentes a 12 coletas não
    usadas no treinamento. Os hiperparâmetros do Extra Trees e do Random Forest
    foram selecionados exclusivamente no treino por validação cruzada estratificada
    em quatro partições, agrupada por coleta. SVM e MLP foram precedidos por
    padronização dos atributos. Para o Extra Trees, foram usados diretamente os
    resultados da matriz de confusão do melhor modelo: 9.377 verdadeiros da classe
    com nematoide, 1.282 erros dessa classe, 10.710 verdadeiros da classe saudável
    e 746 erros dessa classe.}
\end{table}
